% ===================================================================================
% PRÓLOGO
% ===================================================================================
\addchap{Prólogo: Más Allá de la Ecuación, el Valor}
\label{chap:prologo}

Este libro nació de una frustración. O, para ser más precisos, de una brecha.\\

Durante años, ya sea en laboratorios universitarios o en proyectos de consultoría, vi un patrón repetirse: el mundo académico domina el arte de derivar ecuaciones con una elegancia impecable, mientras que el mundo industrial necesita resolver problemas con una urgencia brutal. Entre ambos extremos hay un abismo. De un lado, \textit{papers} que terminan con un ``los resultados son prometedores''; del otro, ingenieros de proyecto que necesitan saber, para el viernes, si un diseño se va a sobrecalentar, deformar o fracturar.\\

¿Dónde estaba el manual que tradujera el poder de los ``primeros principios'' físicos en decisiones de ingeniería concretas, robustas y accesibles?\\

\textit{El Código de la Realidad} es mi intento de escribir ese manual. Es el puente entre la elegante abstracción de las matemáticas y la realidad, a veces sucia y siempre compleja, de la industria. Especialmente de la industria chilena, con sus gigantescas palas mecánicas, sus molinos que trituran montañas, sus sistemas de energía que desafían el desierto y el frío, y su necesidad imperiosa de hacer más con menos: menos energía, menos agua, menos mantenimiento, menos riesgo.\\

Este no es un libro para quedarse en la estantería. Está diseñado para ser usado en el campo de batalla del diseño y la operación industrial. Su filosofía es simple, y se la robé descaradamente al gran Richard Feynman: ``Lo que no puedo crear, no lo entiendo''. Por eso, cada capítulo es una travesía que parte de la ley física fundamental (la ``creación'' de la ecuación) y culmina con su encarnación en código de simulación ejecutable (la ``comprensión'' práctica). Aprenderá no solo la teoría detrás de la ecuación de calor, sino a programar en \fenics la simulación que le dirá exactamente cuánto tarda en fallar un componente crítico.\\

\textbf{A quién va dirigido este libro:}
\begin{itemize}
    \item Al físico o ingeniero recién graduado que quiere que su profundo conocimiento teórico impacte en el mundo real.
    \item Al analista o ingeniero de proyecto en la industria que intuye que la simulación puede darle una ventaja, pero no sabe por dónde empezar.
    \item Al estudiante avanzado que está cansado de los ejercicios de juguete y anhela enfrentarse a problemas con el aroma de la realidad.
    \item Al futuro miembro del equipo de Wasaff Consulting, para quien este libro será el punto de partida de un riguroso entrenamiento.
\end{itemize}

\textbf{Qué encontrará dentro:}
Un viaje desde lo simple hasta lo complejo. Comenzaremos con los pilares de toda simulación (el equilibrio y la evolución temporal) para luego adentrarnos en el fascinante mundo de los problemas no lineales y, finalmente, destrabar el poder de la multifísica: cómo el calor deforma los sólidos, cómo los fluidos transportan energía y cómo la turbulencia y la interacción fluido-estructura dictan el comportamiento de los sistemas más complejos. Cerraremos el círculo mostrando cómo usar estas simulaciones no solo para predecir, sino para optimizar y diseñar, que es, al final, el verdadero objetivo.\\

Chile tiene todos los desafíos y el talento para enfrentarlos. Mi esperanza es que este libro sirva como una herramienta para que físicos, ingenieros y empresas chilenas no solo entiendan el mundo, sino que lo modelen, lo simulen y, en última instancia, lo rediseñen para ser más eficiente, seguro y sostenible.\\

El código de la realidad está escrito en lenguaje matemático. Estemos listos para leerlo.

\vspace{2cm} % Añade un espacio vertical antes de la firma

\begin{flushright}
    \textit{Felipe Wasaff} \\
    Fundador de Wasaff Consulting \\
    Septiembre de 2025
\end{flushright}

\vspace{2cm}

\begin{center}
    \textsc{Wasaff Consulting} \\
    \url{www.wasaffconsulting.com} \\
    \textit{De la Física a la Solución.}
\end{center}
